% Chapter 1
\chapter{Introduction to Shells} % Main chapter title

\label{Chapter1} % For referencing the chapter elsewhere, use \ref{Chapter1} 

%----------------------------------------------------------------------------------------

% Define some commands to keep the formatting separated from the content 
\newcommand{\keyword}[1]{\textbf{#1}}
\newcommand{\tabhead}[1]{\textbf{#1}}
\newcommand{\code}[1]{\texttt{#1}}
\newcommand{\file}[1]{\texttt{\bfseries#1}}
\newcommand{\option}[1]{\texttt{\itshape#1}}

%----------------------------------------------------------------------------------------
\section{What are shells and membranes?}
Most structural objects and elements can be reduced to a physical line with a certain cross-section. The description of stresses
 in these objects can be done by describing  “stress resultants” - the normal force, transverse shearing forces, bending
 moments, and the torque - in the cross section. In contrast to these objects, certain structural elements enclose some space making it impossible for them to be reduced to a line. They are better described as a curved surface
 or a plane. These elements can be broadly broken down in two categories - plates and shells \cite{flugge2013}.

A shell is a curved, thin-walled structure with the thickness being significantly smaller than the other dimensions \cite{bischoff2004}.
 The distinguishing feature between a shell and a plate is that the shell has a curvature whereas plates are flat
 structures \cite{chen1997}. It is, however, not necessary for the thickness in a shell to be extremely small. It can also vary across
 the surface. The behaviour of a shell can be geometrically linear or nonlinear. Membranes are special cases of shells whose walls offer
 no resistance to bending moments and transverse shear forces. 

Shells carry load efficiently - this is most optimally done when all the load is carried using membrane forces; 
 exhibit high structural integrity, strength to weight ratio; and have high stiffness \cite{ventsel2001}.
 Shells and membranes are ubiquitous in engineering and nature because of the advantages they present - from aerospace as
 aeroplanes and rocket hulls, civil engineering where they are found in large span roofs, cooling towers, containment
 tanks, to biological systems such as skulls, eyes, cells, and blood vessels.
%----------------------------------------------------------------------------------------
%----------------------------------------------------------------------------------------
\section{Problems with modelling shells/membranes}
In the real world, thin shells are very sensitive to imperfections and can be highly unstable if not designed properly \cite{loibl2019}.
 Modelling this complexity is challenging too. Shells can undergo large deformations, large rotations, and large strains.
 They can behave highly nonlinearly and can have multiple instabilities - in compression, for out-of-plane loading, for 
 pressure loading, and with respect to in-plane loading \cite{sauer2014}.
 This mixture of performance and sensitivity makes the shell the 'primadonna of structures' \cite{bischoff2004}. Given that there are
 no analytical solutions for most cases to predict a shell's behaviour, reliable approximate solutions are necessary.

Shell theories aim to reduce the dimensionality of the equilibrium problem from a 3D body to only the study of the
 midsurface of the shell. This dimensionality reduction is often dependent on the thickness of shells. Generally, shells
 with a thickness h and general radius of curvature R, can be categorised as thin shells if $\text{max}(h/R)<=(1/20)$.
Love laid the foundations of classical linear shell theory by applying the Kirchhoff hypothesis to thin,
 homogenous, isotropic, linear-elastic shells governed by a generalised Hooke's law. Known as the Kirchhoff-Love theory,
 the assumptions can be summarised as \cite{podio1989}:
\begin{itemize}
    \item No deformation in the midsurface - it is the neutral plane during bending.
    \item Points that lie on the normal to the midsurface remain so even after deformation.
\end{itemize}
The transverse normal stresses and the transverse shear strains are neglected. 
Reissner \cite{reissner1941} and Mindlin \cite{mindlin1951} separately expanded upon the Kirchhoff-Love hypothesis with the inclusion of transverse shear
strains and first order shear effects in their derivations, which could calculate the deformation of shells with a
thickness of $\text{max}(h/R) <= (1/10)$. Their theories assumed the following:
\begin{itemize}
    \item The normals of the undeformed midsurface are deformed into the normals to the deformed midsurface.
    \item The stress components normal to the middle surface are small compared to the other stress components and can be neglected in the stress-strain relations
    \item The displacements are so small that the equilibrium conditions for deformed elements are the same as if they were not deformed.
\end{itemize}
There is a lack of a universally accepted classical linear shell theory, all of which share a degree of inaccuracy
 due to their arbitrary treatment of secondary effects \cite{nayfeh2004} and \emph{ad hoc} assumptions \cite{naghdi1963}
 during different stages of the formulation of the governing equations. Love's first
 approximation
 \begin{itemize}
    \item doesn't satisfy the equilibrium equation arising from the symmery of the stress tensor,
    \item isn't invariant under infinitesimal rigid body displacement and,
    \item violates a 2D analogue of the Betti-Rayleigh reciprocity theorem in 3D linear elasticity.
\end{itemize}
Reissner's first approximation runs into the same problems except violating the reciprocity theorem \cite{truesdell1984}. Although improvements to classical linear
 shell theories can be found in later works by Reissner (1959), Novozhilov (1959), Sanders (1959), and Koiter
 (1959), Naghdi \cite{naghdi1963} mentions the existence of inconsistencies present in many of these approximations. 

The formulation presented in this report is based on the work of Sauer \cite{sauer2014}, which is a formulation of
 the theoretical framework laid down by Steigmann \cite{steigmann1999}. Although substantially same as Naghdi's \cite{truesdell1984}
 formulation, Steigmann's assumptions logically conform to finite elasticity theory. Sauer's formulation is also specifically based on
 membranes because while membranes are special cases of shells and a formulation for shells could be extended to membranes, a formulation dedicated explicitly
 towards membranes is more robust and accurate \cite{sauer2014}. Here, liquid membranes can also be accounted for through the
 split of contributions between the in-plane and out-of-plane components, admission of arbitrary elastic material models,
 all quantities are formulated directly in curvilinear coordinates.

The following chapter presents the theory and formulation of the finite element method for the solution of an elliptic
 boundary vaue problem. Chapter 3 details the theory and implementation of the finite element method to nonlinear boundary
 value problems for solid mechanics in an orthonormal basis system. Chapter 4 extends this nonlinear framework to membranes
 through the incorporation of curvilinear coordinates and presents numerical examples to illustrate the capabilities of the
 formulation.
%----------------------------------------------------------------------------------------
\afterpage{\blankpage}
